% 题干转写：完成
% 解析转写：完成
% 图形核对：完成
\section{第35届复赛第3题：细杆打击中心与撤轴后运动}
\subsection{题干}
质量为 $M$、长度为 $l$ 的匀质细杆 $AB$ 自由悬挂于通过原点 $O$ 的水平光滑转轴上，初始时杆竖直下垂。质量为 $m$ 的弹丸以速度 $v_0$ 水平撞击杆的打击中心并嵌入杆中；杆转过半圈恢复竖直状态时立即撤去转轴。忽略地面高度限制，重力加速度为 $g$。如图\ref{fig:cpho35q03}，求打击中心、撤轴前轴力、撤轴后杆端坐标，以及杆再转半圈时 $O,B$ 两点的高度差。
\begin{figure}[htbp]\centering
\begin{tikzpicture}[>=Stealth,scale=0.82]
  \begin{scope}
    \draw[->,thick] (-2.3,0) -- (2.3,0) node[right] {$x$};
    \draw[->,thick] (0,-2.9) -- (0,1.0) node[above] {$y$};
    \fill (0,0) circle (1.5pt) node[above left] {$O$};
    \draw[very thick] (0,0) -- (0,-2.35);
    \node[left] at (-0.10,-1.2) {$M,l$};
    \fill (0,-2.35) circle (1.3pt) node[right] {$B$};
    \draw[->,very thick,red!75!black] (2.0,-1.58) -- (0.72,-1.58) node[midway,above] {$v_0$};
    \fill[blue!65] (0.60,-1.58) circle (3pt) node[right] {$m$};
    \draw[densely dashed] (0,-1.58) -- (0.55,-1.58);
    \node[left] at (-0.1,-1.58) {$L=2l/3$};
    \node[below] at (0,-3.12) {初始撞击状态};
  \end{scope}
  \begin{scope}[xshift=6.0cm]
    \draw[->,thick] (0,0) -- (2.7,0) node[right] {$x'$};
    \draw[->,thick] (0,0) -- (0,2.2) node[above] {$y'$};
    \fill (0,0) circle (1.4pt) node[below left] {质心};
    \draw[very thick] (0,0) -- (1.35,1.35);
    \fill (1.35,1.35) circle (1.2pt) node[above right] {$B$};
    \draw[->,thick,purple!75!black] (0.48,0.48) arc[start angle=45,end angle=-20,radius=0.68];
    \node[purple!75!black] at (0.93,0.15) {$\omega't$};
    \node[below] at (0,-0.55) {撤轴后相对质心转动};
  \end{scope}
\end{tikzpicture}

\caption{第35届复赛第3题题图与撤轴后相对质心坐标的TikZ重绘。}\label{fig:cpho35q03}
\end{figure}

\subsection{参考解答整理}
\paragraph{（1）打击中心。} 设作用点到 $O$ 的距离为 $L$。由杆绕端点的转动方程与质心平动关系，
\[FL=\frac13Ml^2\beta,\qquad F-F_x=Ma_x,\qquad a_x=\frac l2\beta.\]
打击中心要求撞击瞬间 $F_x=0$，故
\[\boxed{L=\frac23l.}\]

\paragraph{（2）撤轴前的轴力。} 碰撞后杆--弹丸系统的质心到轴距离、初始角速度和转动惯量为
\[
L_c=\frac{4m+3M}{6(M+m)}l,\qquad
\omega_0=\frac{6mv_0}{(4m+3M)l},\qquad
I=\frac19(3M+4m)l^2.
\]
转过 $\theta$ 后，机械能守恒给出
\[
\omega^2=\frac{36m^2v_0^2}{(4m+3M)^2l^2}-\frac{3g}{l}(1-\cos\theta),
\qquad \beta'=\frac{3g}{2l}\sin\theta.
\]
轴对系统作用力的分量为
\[
N_x=\frac{(4m+3M)g}{4}(3\sin\theta\cos\theta-2\sin\theta)
+\frac{6m^2v_0^2}{(4m+3M)l}\sin\theta,
\]
\[
N_y=(M+m)g+\frac{6m^2v_0^2}{(4m+3M)l}\cos\theta
+\frac{(4m+3M)g}{4}(3\cos^2\theta-2\cos\theta-1).
\]

\paragraph{（3）撤轴后的杆端轨迹。} 在 $\theta=\pi$ 时撤轴，角速度记为
\[
\omega'=\sqrt{\frac{36m^2v_0^2}{(4m+3M)^2l^2}-\frac{6g}{l}}.
\]
此后质心作平抛、杆绕质心顺时针转动。以撤轴瞬间为 $t=0$，杆端 $B$ 的坐标为
\[
x_B(t)=L_c\omega't+(l-L_c)\sin(\omega't),
\]
\[
y_B(t)=L_c-\frac12gt^2+(l-L_c)\cos(\omega't).
\]

\paragraph{（4）再转半圈时的高度差。} 令 $\omega't_1=\pi$，则
\[
h=\left|2L_c-l-\frac12g\left(\frac\pi{\omega'}\right)^2\right|
=\left|\frac{\pi^2g}{2}\left[\frac{36m^2v_0^2}{(4m+3M)^2l^2}-\frac{6g}{l}\right]^{-1}-\frac{ml}{3(M+m)}\right|.
\]

\subsection{转写依据}
题干、答案公式和题图分别依据 \texttt{CPHO35-Q03.md}、\texttt{CPHO35-Q03\_formula\_check\_v1.csv} 及图件核对入口；当前PDF未明示分值，未写入图册。
